\documentclass[12pt]{article}
\usepackage[utf8]{inputenc}
\usepackage[T1]{fontenc}
\usepackage{hyperref}
\title{Bathari Language Resources for K-12 Curriculum Development}
\author{Compiled by Claude Sonnet 5.5 (claude-sonnet-5-5) for Andr\'as Kornai}
\date{}
\begin{document}
\maketitle

\section{Introduction}
Bathari (Ba\d{t}\d{h}ari, ISO 639-3 \texttt{bhm}) is a Modern South Arabian (Semitic, Afro-Asiatic) language of Dhofar, Oman. It is critically endangered; estimates of fluent speakers range from about a dozen to a few dozen, mostly elderly, and some members of the Bathari tribe speak Mehri instead. Very little material exists, and none that is designed for schools.

\section{Primary Educational Resources}
No primers, textbooks or school programmes for Bathari were found in the search results.

\section{Linguistic and Reference Materials}
\begin{itemize}
\item Wikipedia overview of the language: \url{https://en.wikipedia.org/wiki/Ba%E1%B9%AD%E1%B8%A5ari_language}
\item DICE (University of Missouri) Afro-Asiatic questionnaire document on Bathari: \url{https://dice.missouri.edu/assets/docs/afro-asiatic/Bathari.pdf}
\end{itemize}
A descriptive grammar by Fabio Gasparini (``The Ba\d{t}\d{h}ari Language of Oman -- Towards a Descriptive Grammar'') was mentioned in search results, but no URL was seen and it was not examined, so it is not listed as a resource here.

\section{Religious and Community Resources}
None found.

\section{Digital and Audio Resources}
\begin{itemize}
\item Endangered Languages Project entry: \url{https://endangeredlanguages.com/lang/1595}
\end{itemize}

\section{Curriculum Development Context}
With so few speakers and almost no published teaching material, any curriculum would depend on new documentation work with the remaining speakers, building on the existing descriptive scholarship. Closely related Mehri and Jibbali materials may help, but were not searched for here.
\end{document}